% GENERATED FILE. Edit canonical lessons, then run npm run generate:books.
% canonical-source-hash: fnv1a64:f08ef63dbb6d826f
% canonical-lessons: FR-C284-reveler, FR-C284-rapporter, FR-C284-dementir, FR-C284-relater, FR-C284-commenter

\chapter{Ce que disent les journaux}
\label{ch:fr-ce-que-disent-les-journaux}
\hlchaptermodality{284}

\begin{chapteropening}
\textbf{By the end of this chapter:} \emph{I can report what people said and found out: revealing, reporting, denying, recounting and commenting.}

Complete the last lesson of chapter 284: I can report what people said and found out: revealing, reporting, denying, recounting and commenting.
\end{chapteropening}

\section[révéler]{révéler — to reveal, to disclose}

\label{lesson:FR-C284-reveler}



\begin{quote}
Before the new one: say the French for to interview, then the French for the aerial.
\end{quote}

\subsection*{You'll want to know: révéler}
\textbf{révéler} — \textquotedblleft{}to reveal, to disclose\textquotedblright{}.

\begin{quote}
Le journal a révélé le secret du chanteur.
\end{quote}

\begin{quote}
\emph{The paper revealed the singer's secret.}
\end{quote}

\emph{Se révéler} — to turn out: \emph{il s'est révélé très utile}.

\begin{cousinweb}
\emph{révéler} is a learned borrowing from Latin \emph{revēlāre} \textquotedblleft{}to unveil\textquotedblright{} (EtymDB), from \emph{re-} \textquotedblleft{}back\textquotedblright{} and \emph{vēlum} \textquotedblleft{}a veil\textquotedblright{}.
\end{cousinweb}

\begin{grammarlens}[title={What you've built}]
One more word for telling, explaining and describing at B1.
\end{grammarlens}

\subsection*{Your turn}
\begin{itemize}
\raggedright
  \item \emph{Say it:} \emph{révéler}
  \item \emph{Say it:} \emph{révéler}, once more
  \item \emph{Say it:} a sentence of your own with \emph{révéler}
  \item \emph{Recall:} say the French for to interview, then the French for the aerial, then say \emph{révéler} again
\end{itemize}

\begin{tcolorbox}[breakable,colback=teal!4,colframe=teal!35!black,title={Before you move on}]
What does révéler mean? (\textquotedblleft{}To reveal, to disclose\textquotedblright{}.) Say it once more.
\end{tcolorbox}
\section[rapporter]{rapporter — to report, to relate; to bring back}

\label{lesson:FR-C284-rapporter}



\begin{quote}
Before the new one: say the French for the aerial, then the French for to reveal.
\end{quote}

\subsection*{You'll want to know: rapporter}
\textbf{rapporter} — \textquotedblleft{}to report, to relate; to bring back\textquotedblright{}.

\begin{quote}
Le journaliste a rapporté ce qu'il avait vu.
\end{quote}

\begin{quote}
\emph{The journalist reported what he had seen.}
\end{quote}

Also to bring back: \emph{rapporter des souvenirs de vacances}.

\begin{cousinweb}
\emph{rapporter} is built on \emph{apporter} \textquotedblleft{}to bring\textquotedblright{}, which comes from Latin \emph{apportāre} (EtymDB), from \emph{ad} and \emph{portāre} \textquotedblleft{}to carry\textquotedblright{} (Lewis \& Short).
\end{cousinweb}

\begin{grammarlens}[title={What you've built}]
One more word for telling, explaining and describing at B1.
\end{grammarlens}

\subsection*{Your turn}
\begin{itemize}
\raggedright
  \item \emph{Say it:} \emph{rapporter}
  \item \emph{Say it:} \emph{rapporter}, once more
  \item \emph{Say it:} a sentence of your own with \emph{rapporter}
  \item \emph{Recall:} say the French for the aerial, then the French for to reveal, then say \emph{rapporter} again
\end{itemize}

\begin{tcolorbox}[breakable,colback=teal!4,colframe=teal!35!black,title={Before you move on}]
What does rapporter mean? (\textquotedblleft{}To report, to relate; to bring back\textquotedblright{}.) Say it once more.
\end{tcolorbox}
\section[démentir]{démentir — to deny, to refute (officially)}

\label{lesson:FR-C284-dementir}



\begin{quote}
Before the new one: say the French for to reveal, then the French for to report.
\end{quote}

\subsection*{You'll want to know: démentir}
\textbf{démentir} — \textquotedblleft{}to deny, to refute (officially)\textquotedblright{}.

\begin{quote}
Le chanteur a démenti la nouvelle.
\end{quote}

\begin{quote}
\emph{The singer denied the news.}
\end{quote}

Conjugated like \emph{sortir}: \emph{il dément, ils démentent}.

\begin{cousinweb}
\emph{démentir} is built on \emph{dé-} and the verb for \textquotedblleft{}to lie\textquotedblright{} (EtymDB), which is inherited from Latin \emph{mentīrī} \textquotedblleft{}to lie\textquotedblright{} (Lewis \& Short): to show that something is false.
\end{cousinweb}

\begin{grammarlens}[title={What you've built}]
One more word for telling, explaining and describing at B1.
\end{grammarlens}

\subsection*{Your turn}
\begin{itemize}
\raggedright
  \item \emph{Say it:} \emph{démentir}
  \item \emph{Say it:} \emph{démentir}, once more
  \item \emph{Say it:} a sentence of your own with \emph{démentir}
  \item \emph{Recall:} say the French for to reveal, then the French for to report, then say \emph{démentir} again
\end{itemize}

\begin{tcolorbox}[breakable,colback=teal!4,colframe=teal!35!black,title={Before you move on}]
What does démentir mean? (\textquotedblleft{}To deny, to refute (officially)\textquotedblright{}.) Say it once more.
\end{tcolorbox}
\section[relater]{relater — to recount, to relate}

\label{lesson:FR-C284-relater}



\begin{quote}
Before the new one: say the French for to report, then the French for to deny.
\end{quote}

\subsection*{You'll want to know: relater}
\textbf{relater} — \textquotedblleft{}to recount, to relate\textquotedblright{}.

\begin{quote}
Le livre relate la vie d'une famille pendant la guerre.
\end{quote}

\begin{quote}
\emph{The book recounts the life of a family during the war.}
\end{quote}

Formal and written; in speech, \emph{raconter}.

\begin{cousinweb}
\emph{relater} is built on Latin \emph{relātus} (EtymDB), the past participle of \emph{referre} \textquotedblleft{}to carry back, to report\textquotedblright{} (Lewis \& Short).
\end{cousinweb}

\begin{grammarlens}[title={What you've built}]
One more word for telling, explaining and describing at B1.
\end{grammarlens}

\subsection*{Your turn}
\begin{itemize}
\raggedright
  \item \emph{Say it:} \emph{relater}
  \item \emph{Say it:} \emph{relater}, once more
  \item \emph{Say it:} a sentence of your own with \emph{relater}
  \item \emph{Recall:} say the French for to report, then the French for to deny, then say \emph{relater} again
\end{itemize}

\begin{tcolorbox}[breakable,colback=teal!4,colframe=teal!35!black,title={Before you move on}]
What does relater mean? (\textquotedblleft{}To recount, to relate\textquotedblright{}.) Say it once more.
\end{tcolorbox}
\section[commenter]{commenter — to comment on}

\label{lesson:FR-C284-commenter}



\begin{quote}
Before the new one: say the French for to deny, then the French for to recount.
\end{quote}

\subsection*{You'll want to know: commenter}
\textbf{commenter} — \textquotedblleft{}to comment on\textquotedblright{}.

\begin{quote}
Il commente les matchs à la radio.
\end{quote}

\begin{quote}
\emph{He commentates on matches on the radio.}
\end{quote}

No \emph{sur} after it: \emph{commenter un match}; with the noun, \emph{faire un commentaire sur}.

\begin{cousinweb}
\emph{commenter} is a learned borrowing from Latin \emph{commentārī} \textquotedblleft{}to think over, to write notes on\textquotedblright{} (Lewis \& Short), which is linked to the root of \emph{mēns} \textquotedblleft{}mind\textquotedblright{}.
\end{cousinweb}

\begin{grammarlens}[title={What you've built}]
One more word for telling, explaining and describing at B1.
\end{grammarlens}

\subsection*{Your turn}
\begin{itemize}
\raggedright
  \item \emph{Say it:} \emph{commenter}
  \item \emph{Say it:} \emph{commenter}, once more
  \item \emph{Say it:} a sentence of your own with \emph{commenter}
  \item \emph{Recall:} say the French for to deny, then the French for to recount, then say \emph{commenter} again
\end{itemize}

\begin{tcolorbox}[breakable,colback=teal!4,colframe=teal!35!black,title={Before you move on}]
What does commenter mean? (\textquotedblleft{}To comment on\textquotedblright{}.) Say it once more.
\end{tcolorbox}
